\section{Background and Threat Model}
\label{sec:background}

\subsection{Honeypots and the Fidelity--Safety Tension}
\label{sec:tension}
A honeypot is a resource with no production value whose only purpose is to
be probed, attacked, and compromised~\cite{spitzner2002}. The literature
classifies honeypots by interaction depth~\cite{lopez2024}. Low-interaction
systems emulate a handful of protocols well enough to be scanned;
medium-interaction systems accept sessions and expose a bounded command
surface; high-interaction systems hand the attacker a real machine, with
the containment cost that implies. Production deployments overwhelmingly
favor low- and medium-interaction designs because they contain risk:
a scripted trap cannot be pivoted through, and its worst-case disclosure
is the script itself. Containment, however, is bought by scripting, and
the deception is a finite automaton. A finite automaton cannot answer a
question it was not compiled to answer, so an attacker who strays off the
script, or simply probes from a class the author did not anticipate,
learns nothing new and may learn the host is fake.

Deception research gives the tension a formal shape: a deception succeeds
when the attacker's belief about the system diverges, in a way the
operator benefits from, from the system's reality~\cite{bera2024}, and
models of attacker--defender interaction treat the attacker's
discrimination between real and fake as the quantity the deception must
suppress~\cite{shinde2020,liska2024}. That framing exposes the catch-22
that confines classical designs. A honeypot is valuable exactly insofar
as the attacker believes it and engages with it, and engagement rewards
richness: more services, more answers, more history. Richness, delivered
by scripts, is where the risk lives, because every scripted answer is a
hand-written commitment about a machine that does not exist, and an
attacker who walks off the script collects a contradiction or a crash
instead of a session. Classical practice manages this by keeping the
fiction poor. The interaction-depth taxonomy states the trade as
fidelity against safety; stated in the adversary's terms, it is a
catch-22: the fiction must be rich to be believed and poor to be safe.

Raising interaction depth through generative models inverts the trade.
A generator can answer \emph{anything}, which is exactly the problem.
First-generation generative honeypots demonstrate that an LLM can sustain
a shell conversation, an LDAP directory, or an AI-agent lure convincingly
enough to hold modern attackers~\cite{valeros2024,reworr2024,jimenez2025};
their evaluations concentrate on believability and engagement. What the
emerging line leaves unexamined is the complementary question: what
happens when the model is \emph{wrong}? For a security tool, wrongness
takes three forms, and all three are silent failures. First,
\emph{inconsistency}: the fiction drifts, and an attacker who
cross-references two turns learns that the host is synthetic. Second,
\emph{over-disclosure}: the model complies too well, emits credential
material that was never intended to be served, or produces content that
poisons downstream analytics. Third, \emph{unavailability}: the generator
is a rented service that can refuse, stall, or vanish mid-session, and a
deception endpoint that answers with a refusal has told the attacker
something true. The design in \S\ref{sec:design} treats all three as
first-class constraints with enforcing mechanisms, not quality hopes.
Resolving the catch-22, rather than trading within it, is the design
problem this paper addresses: \S\ref{sec:design} decouples the realism
of the fiction from the risk of the machinery that produces it, by
making every generator output pass a check computed against a world the
defender fully controls.

\subsection{The Intelligence Ladder: Trust Prices and Exit Ramps}
\label{sec:rungs}

The design question this paper asks is how much intelligence a serving
loop should spend. We answer it with a pricing rule: a capability is
admitted to the serving path only if its measured marginal benefit
(\S\ref{sec:evaluation}) justifies its \emph{trust price} --- the
attack surface it opens, the certification burden it places on
everything downstream of it, and the availability dependency it
introduces. The loop answers the spending question with a single axis,
\emph{certifiability} --- whether a mechanism's output can be checked
against the closed world before it reaches the wire --- and spends
along it in rungs that compose: the ladder of \S\ref{sec:ladder}.
Table~\ref{tab:rungs} prices the four rungs the system implements and
declares each rung's \emph{exit ramp}: the deterministic behavior the
loop takes when the rung fails, so that no failure mode of any rung is
left to improvisation.

\begin{table}[t]
\centering
\caption{The four rungs of the intelligence ladder: what each
delegates to a less certifiable mechanism, its trust price, and its
declared exit ramp.}
\label{tab:rungs}
\scriptsize
\setlength{\tabcolsep}{2.5pt}
\begin{tabular}{p{0.13\columnwidth}p{0.31\columnwidth}p{0.45\columnwidth}}
\toprule
Rung & Decision delegated & Trust price and exit ramp \\
\midrule
R0 templates & none: serving itself & No intelligence to price; the floor that every other rung's exit ramp lands on (\S\ref{sec:fallback}) \\
R1 classifier & which deception policy a session receives & Cheap, deterministic at inference, auditable offline; a misroute changes what is served, never whether it is verified, and the deterministic behavior floor subsumes adaptive rerouting at $\Delta{=}0$ (\S\ref{sec:jeveval}) \\
R2 retrieval planner & which intelligence shapes the plan & Statistics over an open corpus upstream of the gate; failed or refused retrieval degrades the plan, never the serving rule (\S\ref{sec:don}) \\
R3 generator & what text to serve & The only rung not certifiable in advance; exit: one violation-fed retry, then R0, with \textsc{require} content stopping at the human gate (\S\ref{sec:verify}, \S\ref{sec:hitl}) \\
\bottomrule
\end{tabular}
\end{table}

The rule cuts in both directions: intelligence is admitted at a known
price, and a rung whose measured benefit can be reproduced by a
cheaper, certifiable mechanism should be compiled down.
\S\ref{sec:design} gives the mechanism and the descent;
\S\ref{sec:evaluation} gives each rung's measured marginal.

\subsection{Threat Model}
\label{sec:threat-model}
\label{sec:threatmodel}
We consider an attacker who interacts with the honeypot over standard
remote-access protocols (SSH today, HTTP and Telnet supported by the same
machinery). The attacker may be unskilled (credential spraying, banner
grabbing) or expert (hands-on-keyboard, staged tooling, lateral-movement
reconnaissance); the behavioral difference between the two is what the
profiler must recover. The attacker's goals are those of the CTI
literature: initial access, credential harvest, privilege escalation, and
exfiltration, observable as techniques in the \attck matrix. Sessions are
attributed behaviorally, not cryptographically: the system may link
reconnaissance in one session to credential use in another, but an
attacker who rotates infrastructure presents as a new profile.

The attacker's \emph{observation model} is the served response and
nothing else. They cannot see the plan that produced it, the graph that
constrained it, or the verdicts that preceded it. Within that constraint
we assume arbitrary fingerprinting skill: the attacker may issue
conflicting commands, re-query facts stated many turns earlier, compare
notes across parallel sessions, and probe protocol edge cases
specifically to test whether the host is synthetic. Inconsistency is
therefore not a cosmetic defect but an information channel to the
adversary, which is why G2 (below) requires a zero-contradiction
condition rather than a low rate.

Timing is the one channel this observation model leaks. The attacker
measures only responses, but response latency is measurable: served
responses average ${\sim}44$~s in the live campaign, dominated by
generative turns (\S\ref{sec:discussion}), producing a shell
interaction that feels unlike a local host; the
fallback's fast vetted answers have a signature of their own. We treat
latency homogenization as an open problem rather than a solved one, but
not an unmeasured one: \S\ref{sec:timing-measured} reports an
attacker-seat ROC analysis of this channel (unmitigated AUC $0.983$; a
timing-shaped deterministic path reaches the real-vs-real drift floor)
and the deception properties this paper claims remain those verifiable
from response content alone.

We do \emph{not} assume the attacker can compromise the honeypot's
control plane, escape the container sandbox, observe traffic between
honeypot components, or influence the retrieval corpus before a session
begins. Physical access, insider threat, and compromise of the defender's
workstation are out of scope.

Finally, we assume the attacker may be \emph{adaptive}: they may probe
systematically to map the fiction's boundaries, deliberately elicit
refusals to fingerprint the generator, re-attack across sessions to test
which facts persist, and seed inputs shaped to influence what the
retrieval engine surfaces. The design does not assume the deception
logic is secret; the fallback transcript library is assumed
disclosable, and the fiction's single-scope deployment means an attacker
who compromises one session's view learns one host, not the fleet. What
the system does rely on keeping secret is small: the canary registry,
which must exist outside any view the attacker can reach, and the
serving-time verdicts of \S\ref{sec:design}. An adaptive adversary is
the reason the evaluation of \S\ref{sec:evaluation} measures the system
under targeted adversarial re-query rather than benign traffic.

We additionally consider the \emph{provider} as an involuntary third
party. The honeypot's generative components run against an inference API
whose safety alignment, rate limits, and availability are outside our
control. Across this project's deployments we observed the full failure
taxonomy firsthand: safety refusals of the deception task, budget
exhaustion under alignment pressure, timeouts, transient network
errors, and catalog churn, in which a listed model disappears from the
provider's offering between one day and the next. A deployment must remain safe and useful
through every row of that taxonomy; \S\ref{sec:design} accordingly treats
the provider as an unreliable component, not a trusted oracle.

Defenders operate the system and are assumed to act within policy. The human-in-the-loop
gate exists precisely because some content classes (credential bait, for
instance) should not be served without a human decision, even when the
machinery would allow it.

\subsection{Vocabulary}
\label{sec:vocab}
Five terms carry precise meanings throughout the paper. A \emph{deception
plan} is the per-session specification the loop serves against: the facts
to project, the bait artifacts to deploy, and the attacker model they are
chosen for. The \emph{fiction graph} is the triple store binding every
projected host fact (hostname, kernel, service versions, credentials) to
one consistent world; \S\ref{sec:design} formalizes its consistency
property. A response is \emph{verified} only if a serving-time check
confirms every factual claim it makes against the graph; anything else
that must be sent is served from the \emph{deterministic fallback}, a
library of human-vetted transcripts. A \emph{honeytoken} is a deployed
artifact (a credential, a configuration value) whose value is a
canary-registered secret, so that any use of it elsewhere is attributable
to the disclosure event. \emph{TTP} denotes an ATT\&CK tactic or
technique~\cite{roy2023}; extraction of TTPs from attacker behavior is
the CTI product the system computes as a by-product of deceiving, and
retrieval-augmented grounding~\cite{lewis2020} over curated CTI corpora,
a technique adapted to security annotation in prior
work~\cite{lekssays2025}, is what supplies that product its evidence base.

\subsection{Design Goals}
\label{sec:goals}
The five goals the design exists to hold --- grounding, consistency,
safety, graceful degradation, and auditability --- are stated with the
mechanisms that enforce and measure them, in \S\ref{sec:design}, so
each goal, its enforcement, and the measurement that can falsify it are
read against the component that carries it.
